% ==================================================================
% id:       comb-029
% topic:    קומבינטוריקה
% subtopic: הסברים קומבינטוריים לזהויות
% solution: full
% ==================================================================
תנו הסבר קומבינטורי לזהויות:
\begin{enumerate}
\item $\displaystyle\sum_{i=0}^n i\binom ni=n\,2^{n-1}$.
\hint{מתוך $n$ סטודנטים בוחרים ועדה (בכל גודל), ובוחרים לה יו"ר מבין חבריה.}
\begin{solbox}
נספור בשתי דרכים את מספר האפשרויות לבחור, מתוך $n$ סטודנטים, ועדה ויו"ר שהוא אחד מחבריה.

\textbf{קודם הוועדה.} אם בוועדה $i$ חברים, בוחרים אותה ב-$\binom ni$ דרכים ואת היו"ר ב-$i$ דרכים. סוכמים על כל הגדלים האפשריים: $\sum_{i=0}^n i\binom ni$. (ועדה ריקה לא תורמת, כי אין לה יו"ר.)

\textbf{קודם היו"ר.} את היו"ר בוחרים ב-$n$ דרכים, ואז כל אחד מ-$n-1$ האחרים בוועדה או לא: $2^{n-1}$ אפשרויות. סך הכול $n\,2^{n-1}$.
\end{solbox}
\item $\displaystyle\binom{2n}2=2\binom n2+n^2$.
\hint{בכיתה $n$ בנים ו-$n$ בנות, ובוחרים שני תלמידים.}
\begin{solbox}
בכיתה $n$ בנים ו-$n$ בנות, ובוחרים שני תלמידים. סך הכול יש $\binom{2n}2$ אפשרויות. מצד שני, נפריד לשלושה מקרים זרים: שני בנים ($\binom n2$), שתי בנות ($\binom n2$), או בן ובת ($n\cdot n=n^2$). לפי עקרון החיבור: $\binom{2n}2=2\binom n2+n^2$.
\end{solbox}
\end{enumerate}
